
%%%%%%%%%%%%%%%%%%%%%%%%%%%%%%%%%%%%%%%%%%%%%%%%%%%%%%%%%%%%%%%
%                                                             %
%             Formules pour l'examen 2 de STT-1000           %
%                                                             %
%%%%%%%%%%%%%%%%%%%%%%%%%%%%%%%%%%%%%%%%%%%%%%%%%%%%%%%%%%%%%%%
\documentclass [12pt]{article}
\renewcommand{\baselinestretch}{1.2}
\usepackage[french]{babel}
\usepackage{amsfonts}
\usepackage{amssymb}
\usepackage{amsmath}
\usepackage{latexsym}
\usepackage{graphics}
\pagenumbering{arabic}
\textheight  23cm
\textwidth 18.5cm \hoffset
-3cm \voffset -2.5cm


\parindent 0pt
\usepackage{graphicx}

\date{}
\begin{document}

\begin{center}
{\bf STT-1000~~~~AIDE-MÉMOIRE POUR L'EXAMEN 2} \\
\end{center}

\noindent
{\large {\bf 1. Estimation ponctuelle}}

\vspace{0.1in}
\noindent
{\bf Biais} Soit $\hat{\theta}$, un estimateur d'un paramètre $\theta$. Le biais de $\hat{\theta}$ est
$$\mbox{Biais}(\hat{\theta})=E(\hat{\theta})-\theta.$$

\noindent
{\bf Estimation de la moyenne} Soit $X_1,\ldots,X_n$ un échantillon aléatoire de taille $n$ d'une population où
$X$ est distribuée selon une loi de moyenne $\mu$. Nous estimons habituellement $\mu$
par la moyenne échantillonnale:
$$\hat{\mu}=\overline{X}=\frac{1}{n}\sum_{i=1}^nX_i.$$

\vspace{0.1in}
\noindent
{\bf Estimation de la variance} Soit $X_1,\ldots,X_n$ un échantillon aléatoire de taille $n$ d'une population où
$X$ est distribuée selon une loi de variance $\sigma^2$. Nous estimons habituellement
$\sigma^2$ par la variance échantillonnale:
$$\hat{\sigma}^2=S^2=\frac{1}{n-1}\sum_{i=1}^n(X_i-\overline{X})^2.$$


\vspace{0.2in}
{\large {\bf 2. Les lois échantillonnales classiques}}

\begin{small}
\begin{displaymath}
{\renewcommand{\arraystretch}{2}
\begin{array}{|c||c|c|c|c|}
\hline
\mbox{Loi}& \mbox{Notation}& \mbox{Moyenne}& \mbox{Variance} &
\mbox{Quantile d'ordre }1-\alpha\\[-0.5cm]
&&&&P\mbox{(Variable$>$Quantile)}=\alpha\\
\hline
\mbox{\small Normale standard} & \displaystyle \mbox{N}(0,1)
& \displaystyle 0
& \displaystyle 1
& \displaystyle Z_{\alpha}
\\[.10cm]
\hline
\mbox{\small Khi-carré avec $k$ degrés de liberté}
& \displaystyle \chi^2_k
& \displaystyle k
& \displaystyle 2k
& \displaystyle \chi^2_{\alpha,k}
\\[.10cm]
\hline
\mbox{\small $t$ de Student avec $k$ degrés de liberté} &
\displaystyle t_k & \displaystyle 0
& \displaystyle \frac{k}{k-2} & \displaystyle t_{\alpha,k}
\\[.10cm]
\hline
\mbox{\small $F$ de Fisher avec $k$ et $m$ degrés de liberté}
& \displaystyle F_{k,m} & \displaystyle \frac{m}{m-2}
& \displaystyle \frac{2 m^2 (k+m-2)}{k(m-2)^2(m-4)}
& \displaystyle F_{\alpha,k,m}\\[.10cm]
\hline
\end{array}}
\end{displaymath}
\end{small}

\vspace{0.05 in}
\begin{enumerate}
\item[$\bullet$] Si $Z \sim N(0,1)$, alors $Z^2 \sim \chi^2_1$. 
\item[$\bullet$] Si $Z_1, \ldots, Z_n$ sont i.i.d. $N(0,1)$, alors $\sum_{i=1}^n Z_i^2 \sim \chi^2_n$.
\item[$\bullet$] Si $Z \sim N(0,1)$, $V \sim \chi^2_v$ et $Z$ et $V$ sont indépendantes, alors $ T = \frac{Z}{\sqrt{V/v}} \sim t_v$.
\end{enumerate}

\newpage
\begin{enumerate}
\item[$\bullet$] Soit $X_1, \ldots X_n$  i.i.d. $N(\mu, \sigma^2)$, alors $$\frac{\overline{X} - \mu}{\sqrt{S^2 / n}} = \frac{\overline{X} - \mu}{S / \sqrt{n}} \sim t_{n-1} \;\;\;\;\;\;\;\;\;\;\;\;\; Y = \frac{(n-1)S^2}{\sigma^2} \sim \chi^2_{n-1}.$$
%\item[$\bullet$] Si $X\sim F_{k,m}$, alors $\displaystyle\frac{1}{X}\sim F_{m,k}$.
\item[$\bullet$] Si $X_1,\ldots,X_n$ sont des variables aléatoires indépendantes
avec $X_i\sim N(\mu_i,\sigma^2_i)$ et si $a_1,\ldots,a_n$ sont des constantes,
alors $$Y=\sum_{i=1}^na_iX_i\sim N(\mu,\sigma^2) \mbox{ avec } \mu=\sum_{i=1}^na_i\mu_i \mbox{ et }\sigma^2=\sum_{i=1}^na_i^2\sigma_i^2.$$
\end{enumerate}


%\vspace{0.2in}
\noindent {\large {\bf 3. Le théorème central limite}}

\vspace{0.1in}
\begin{itemize}
\item
Si $X_1, X_2,..., X_n$ sont des variables aléatoires
indépendantes et identiquement distribuées, avec moyenne $\mu$
et avec variance $0 < \sigma^2 < \infty$ et si $n$ est assez grand
($n\ge  30$), alors
$$
\sum\limits_{j=1}^n X_j\approx N(n\mu, n\sigma^2) \Leftrightarrow
{\sum\limits_{j=1}^n X_j-n\mu\over\sqrt{n\sigma^2}}\approx N(0,1)
\Leftrightarrow
{\overline{X}-\mu\over\sqrt{\sigma^2/n}}\approx N(0,1)
\Leftrightarrow
\overline{X}\approx N(\mu,\sigma^2/n).
$$
\item
Dans le cas où les $X_j$ sont des variables aléatoires
discrètes, on utilise la correction pour la continuité:
\begin{eqnarray*}
{\mathbb P}(X \leq b) &\cong& {\mathbb P}\left(X < b+\frac12\right), \quad {\mathbb
P}(X < b) \cong {\mathbb P}\left(X < b - \frac12\right), \\ {\mathbb P}(X \geq a) &\cong& {\mathbb P}\left(X
> a-\frac12\right), \quad {\mathbb P}(X > a) \cong {\mathbb P}\left(X > a+\frac12\right).
\end{eqnarray*}
\end{itemize}

%\vspace{0.1in}
{\large {\bf 4. Intervalles de confiance} } 

1- Intervalles de confiance \`a $100(1-\alpha)\%$ pour une \textbf{moyenne}:

\begin{table}[h!]
\begin{center}
\begin{tabular}{|c||c|}
\hline
Situation & Intervalle \\[4pt]
\hline
& \\[-12pt]
Loi normale, variance connue & {\large $ \overline{X} \; \; \pm \; \; z_{\alpha/2}\frac{\sigma}{\sqrt{n}}$ }
\\[8pt]
Loi normale, variance inconnue & {\large $\overline{X} \; \; \pm \; \; t_{\alpha/2, n-1}\frac{S}{\sqrt{n}}$}
\\[8pt]
Loi inconnue, $n$ grand & {\large $\overline{X} \; \; \pm \; \; z_{\alpha/2}\frac{S}{\sqrt{n}}$ }\\[8pt]
\hline
\end{tabular}
\end{center}
\end{table}

2- Intervalle de confiance  \`a $100(1-\alpha)\%$ pour la \textbf{variance d'une loi normale}:
$$\left[ \frac{(n-1)S^2}{\chi^2_{\alpha/2, n-1}} \ , \ \frac{(n-1)S^2}{\chi^2_{1-\alpha/2, n-1}}   \right]$$\\

\vspace*{-0.4in}
3-
Intervalle de confiance \`a $100(1-\alpha)\%$ pour une \textbf{proportion}:
$$ \hat{p} \pm z_{\alpha/2} \sqrt{\frac{\hat{p}(1-\hat{p})}{n}}$$

\vspace{0.4in}

4- Intervalles de confiance à 100(1-$\alpha$)\% pour une différence de moyennes: 

\begin{table}[h!]
\begin{center}
\begin{tabular}{|c||c|}
\hline
Situation & Intervalle \\[8pt]
\hline
& \\[-12pt]
Lois normales, variances connues & {\large $ \overline{X}_1 - \overline{X}_2 \pm z_{\alpha_2} \sqrt{ \frac{\sigma^2_1}{n_1} + \frac{\sigma^2_2}{n_2}}$ }
\\[10pt]
Lois normales, variances inconnues, égales & {\large $ \overline{X}_1 - \overline{X}_2 \pm t_{\alpha_2; n_1 + n_2 - 2} S_p \sqrt{ \frac{1}{n_1} + \frac{1}{n_2}}$}
\\[10pt]
Lois normales, variances inconnues, inégales & {\large $ \overline{X}_1 - \overline{X}_2 \pm t_{\alpha_2; \nu} \sqrt{ \frac{S_1^2}{n_1} + \frac{S^2_2}{n_2}} $}\\[10pt]
&  avec {\large $\nu = \displaystyle\frac{(s_1^2/n_1 + s^2_2/n_2)^2}{ \displaystyle \frac{s_1^2/n_1}{n_1 +1} + \frac{s^2_2/n_2}{n_2 + 1} } -2 $ }\\[8pt]
\hline
\end{tabular}
\end{center}
\end{table}

avec $$ S_p^2 = \frac{(n_1-1)S_1^2 + (n_2-1)S_2^2}{n_1 + n_2 - 2} $$

\vspace{0.4in}

5 - Intervalle de confiance à 100(1-$\alpha$)\% pour une différence de proportions: 


$$ \hat{p}_1 - \hat{p}_2  \pm z_{\alpha_2} \sqrt{ \frac{\hat{p}_1(1-\hat{p}_1)}{n_1} + \frac{\hat{p}_2(1-\hat{p}_2)}{n_2}}$$


\vspace{0.2 in}
 {\large\bf 5. Tests d'hypothèses}
\vspace{0.1 in}

\begin{enumerate}
\item[$\bullet$]
$\mbox{ }$
\vspace{-.5cm}

%Les probabilités des erreurs de type I et II sont données par
$\begin{array}{cclll}
\alpha
&=& \mathbb{P}\left(\mbox{erreur de type I}\right)
&=& \mathbb{P}\left(\mbox{Rejeter }H_0~|~H_0\mbox{ est vraie}\right)
\\
\beta &=& \mathbb{P}\left(\mbox{erreur de type II}\right) &=&
\mathbb{P}\left(\mbox{Ne pas rejeter }H_0~|~H_1\mbox{ est
vraie}\right)
\end{array}$\\

\item[$\bullet$] La puissance est égale \`a $1-\beta=\mathbb{P}\left(\mbox{Rejeter }H_0~|~H_1\mbox{ est vraie}\right)$. \\
\item[$\bullet$] Le seuil observé $\alpha^*$ (valeur $p$)
est la probabilité d'obtenir une valeur au moins aussi extr\^eme que celle observée, si on répétait l'expérience et que $H_0$ était vraie. C'est la plus petite valeur de $\alpha$ \`a laquelle on peut rejeter $H_0$.
\end{enumerate}

\vspace*{-0.2 in}
\begin{center}
\begin{tabular}{|c|c|c|}
\hline \multicolumn{3}{|c|}{{\bf Tests sur une moyenne $\mu$}}
\\
\hline &&\\Situation & Statistique du test& $H_{0}:\  \mu =
\mu_{0}$
\\
& sous $H_{0}$  & $H_{1}:\  \mu > \mu_{0}$
\\
& & Rejeter $H_{0}$ si:
\\\hline \hline & &
 \\
$ X_1,...,X_n$ indép. N($\mu$,$\sigma^2$) & $\displaystyle
Z_{0}=\frac{{\overline X} - \mu_{0}}{\sigma /\sqrt{n}} \sim
N(0,1)$ & $Z_{0}>Z_{\alpha}$
\\
où $\sigma^{2}$ est connue &  &
\\
\hline & &
\\
$ X_1,...,X_n$ indép. N($\mu$,$\sigma^2$) & $\displaystyle
t_{0}=\frac{{\overline X} - \mu_{0}}{S / \sqrt{n}}\sim t_{n-1}$ &
$t_{0}>t_{\alpha ,n-1}$
\\
où $\sigma^{2}$ est inconnue& &
\\
\hline
&  &
\\
$ X_1,...,X_n$ indép., loi quelconque, & $\displaystyle
Z_{0}=\frac{{\overline X} - \mu_{0}}{S / \sqrt{n}}\sim N(0,1)$ &
$Z_{0}>Z_{\alpha}$
\\
où $\sigma^{2}$ est inconnue et $n$ est grand &
approximativement &
\\
\hline
\end{tabular}
\end{center}


\begin{center}
\begin{tabular}{|c|c|c|}
\hline \multicolumn{3}{|c|}{{\bf Test sur une variance
$\sigma^2$}}
\\
\hline Situation & Statistique du test & $H_{0}:\  \sigma^2 =
\sigma^2_0$
\\
& sous $H_{0}$
& $H_{1}: \sigma^2 > \sigma^2_0 $
\\
 & & Rejeter $H_{0}$ si:
\\
\hline
\hline
 &  &
 \\
$ X_1,...,X_n$ indép. N($\mu$,$\sigma^2$) & $\displaystyle
\chi^2_0=\frac{(n-1)S^2}{\sigma^2_0}\sim\chi^2_{n-1}$ &
$\chi^2_{0}>\chi^2_{\alpha,n-1}$
\\
$\mu$ est inconnue & &
\\
\hline 
\multicolumn{3}{c}{ }
\\
\hline \multicolumn{3}{|c|}{{\bf Test sur une proportion $p$}}
\\
\hline &&\\Situation & Statistique du test & $H_{0}:\  p = p_{0}$
\\
& sous $H_{0}$ & $H_{1}:\  p > p_{0}$
\\
 & & Rejeter $H_{0}$ si:
\\
\hline \hline
 &  &
 \\
$ X_1,...,X_n$ indép. Bernoulli($p$)
 & $\displaystyle
Z_{0}=\frac{{\hat p} - p_{0}}
{\sqrt{\displaystyle\frac{p_{0}(1-p_{0})}{n}}} \sim N(0,1)$ &
$Z_{0}>Z_{\alpha}$
\\
où $n$ est grand & approximativement &
\\
 & où ${\hat p}=\sum\limits_{i=1}^{n}X_i/n$ &
\\
\hline
\end{tabular}
\end{center}

\newpage
\vspace*{-0.4in}
\begin{center}
\begin{tabular}{|c|c|c|}
\hline
\multicolumn{3}{|c|}{{\bf Comparaison de deux moyennes $\mu_{1}$ et
$\mu_{2}$ (échantillons indépendants)}}
\\
\hline Situation & Statistique du test & $H_{0}:\ \mu_{1} =
\mu_{2}$
\\
& sous $H_{0}$
& $H_{1}:\  \mu_{1} > \mu_{2}$
\\
&
& Rejeter $H_{0}$ si:
\\
\hline
\hline
 $X_{1,1}, ..., X_{1,n_1}$ indép. N($\mu_1$, $\sigma^2_1$) &&\\
 $X_{2,1}, ..., X_{2,n_2}$ indép. N($\mu_2$, $\sigma^2_2$) & $\displaystyle Z_0 =
\frac{\overline X_1 -\overline X_2}{\sqrt{\frac{\sigma_1^2}{n_1} +
\frac{\sigma^2_2}{n_2}}} \sim N(0,1)$ & $Z_0 > Z_{\alpha}$
\\
où $\sigma_{1}^{2}$ et $\sigma_{2}^{2}$ sont connues& &
\\
\hline
$X_{1,1}, ..., X_{1,n_1}$ indép. N($\mu_1$, $\sigma^2_1$)& & \\
$X_{2,1}, ..., X_{2,n_2}$ indép. N($\mu_2$, $\sigma^2_2$)&
$\displaystyle t_0 = \frac{\overline X_1 -\overline X_2}{S_p
\sqrt{\frac{1}{n_1} + \frac{1}{n_2}}} \sim t_{n_1+n_2-2}$ & $t_0 >
t_{\alpha,n_1+n_2-2}$
\\
où $\sigma_{1}^{2}$ et $\sigma_{2}^{2}$ sont inconnues
 & &
\\
mais égales $(\sigma_{1}^{2}=\sigma_{2}^{2})$ & où $S_p =
\sqrt{\displaystyle\frac{(n_1-1)S_1^2 + (n_2-1)S_2^2}{n_1 + n_2
-2}}$ &
\\

\hline
\end{tabular}
\end{center}

\begin{center}
\begin{tabular}{|c|c|c|}
\hline
 \multicolumn{3}{|c|}{{\bf Comparaison de deux moyennes
$\mu_1$ et $\mu_2$ (échantillons appariés)}}
\\
\hline Situation & Statistique du test & $H_{0}:\ \mu_1-\mu_2 = 0$
\\
& sous $H_{0}$ & $H_{1}:\  \mu_1-\mu_2 > 0$
\\
& & Rejeter $H_{0}$ si:
\\
\hline \hline 
$X_{1,1}, ..., X_{1,n}$ indép. N($\mu_1$, $\sigma^2_1$)& & \\
$X_{2,1}, ..., X_{2,n}$ indép. N($\mu_2$, $\sigma^2_2$)
&$\displaystyle t_0 = \frac{\overline D}{S_D/\sqrt n} \sim
t_{n-1}$ & $t_0 > t_{\alpha, n-1}$
\\
$D_j = X_{1j} - X_{2j}$  ({\small $j$=1, ...,
n})&\multicolumn{2}{|c|}{ }
\\
$D_1, ...,D_n$ indép. N($\mu_1-\mu_2$, $\sigma^2_D$) &
\multicolumn{2}{|c|}{$\overline D = \frac{1}{n}
\displaystyle\sum^n_{j=1} Dj~~~,~~~S_D = \sqrt{\frac{1}{n-1}
\displaystyle\sum^n_{j=1} (D_j - \overline D)^2}$}
\\
où $\sigma^2_D$ est inconnue& \multicolumn{2}{c|}{}
\\
\hline \multicolumn{3}{c}{ }
\\

%\hline \multicolumn{3}{|c|}{{\bf Comparaison de deux variances (échantillons indépendants)}}
%\\
%\hline &&\\Situation & Statistique du test & $H_{0}:\  \sigma^2_1 =
%\sigma^2_2$
%\\
%& sous $H_{0}$ & $H_{1}:\  \sigma^2_1 > \sigma^2_2$
%\\
%&  & Rejeter $H_{0}$ si:
%\\
%\hline \hline $X_{1,1}, ..., X_{1,n_1}$ indép. N($\mu_1$, $\sigma^2_1$) & &
%\\
%$X_{2,1}, ..., X_{2,n_2}$ indép. N($\mu_2$, $\sigma^2_2$) & $\displaystyle F_0 = \frac{s_1^2}{s_2^2} {\sim}
%F_{n_1-1, n_2-1}$ & $f_0 > F_{\alpha,n_1-1, n_2-1}$  \\
%où $\sigma_{1}^{2}$ et $\sigma_{2}^{2}$ sont inconnues
% &  & 
%\\
%\hline
\end{tabular}
\end{center}


\begin{center}
\begin{tabular}{|c|c|c|}
\hline \multicolumn{3}{|c|}{{\bf Comparaison de deux proportions
$p_{1}$ et $p_{2}$ (échantillons indépendants)}}
\\
\hline Situation & Statistique du test & $H_{0}:\  p_{1} =
p_{2}$
\\
& sous $H_{0}$ & $H_{1}:\  p_{1} > p_{2}$
\\
&  & Rejeter $H_{0}$ si:
\\
\hline \hline $X_{1}$: nombre de ``succès'' &   &
\\
parmi $n_{1}$ $(n_{1} \geq 30)$ & $\displaystyle Z_0 = \frac{(\hat
p_1 - \hat p_2)} {\sqrt{\displaystyle\hat p (1-\hat p) \left(
\frac{1}{n_1} + \frac{1}{n_2} \right)}} \stackrel{{\rm app}}{\sim}
N(0,1)$ & $Z_0 > Z_{\alpha}$
\\
$X_{2}$: nombre de ``succès'' &  &
\\
parmi $n_{2}$ $(n_{2} \geq 30)$ & où $\displaystyle \hat p =
\frac{(n_1 \hat{p}_1 + n_2 \hat{p}_2)}{n_1 + n_2}$ &
\\
\hline
\end{tabular}
\end{center}

\newpage

{\large\bf 6. Régression linéaire simple}


\vspace{0.5cm}

\textbf{Modèle:}  On observe $(x_1,y_1), (x_2, y_2),..., (x_n,y_n)$ et on
suppose que $y_1, y_2,..., y_n$ sont les réalisations de
variables aléatoires indépendantes $Y_1, Y_2,..., Y_n$ telles
que pour chaque $i$ la loi de $Y_i$ est $N(\beta_0 + \beta_1 x_i,
\sigma^2)$.

\vspace{0.1 in}

\textbf{Quantités utiles:} 

\vspace{-0.4 in}
\begin{align*}
S_{xy}&=\sum\limits_{i=1}^n (x_i-\overline
x)(y_i-\overline y) =  \sum\limits_{i=1}^n x_iy_i - n\;\overline
x\;\overline y \\
S_{xx} &=\sum\limits_{i=1}^n (x_i-\overline x)^2 =
\sum\limits_{i=1}^n x_i^2 - n\;\overline x\;^2\\
S_{yy} &=\sum\limits_{i=1}^n (y_i-\overline y)^2 =
\sum\limits_{i=1}^n y_i^2 - n\;\overline y\;^2 \\
R^2 &= \hat{\beta}_1^2 \frac{S_{xx}}{S_{yy}}
\end{align*}

\textbf{Estimateurs des paramètres:}

$$ \bullet \; \; \hat{\beta}_1 =\frac{S_{xy}}{S_{xx}} \;\;\;\;\;\; \;\;\;\;\;\; \bullet \;\; \hat{\beta}_0 = \overline{y} - \hat{\beta}_1 \overline{x} $$
$$ \bullet \;\; \hat{\sigma}^2 = \frac{1}{n-2} \sum_{i=1}^n (e_i)^2 = \frac{1}{n-2} \sum_{i=1}^n (y_i - \hat{y}_i)^2  = \frac{1}{n-2} \sum_{i=1}^n (y_i - (\hat{\beta}_0 + \hat{\beta}_1 x_i))^2$$

\vspace{0.1in}
\textbf{Test d'hypothèse pour $\beta_1$:}
\vspace{0.1in}

Hypothèse nulle:  $H_0:  \beta_1= \beta_{1,0}$

Statistique du test sous $H_0$ : 
$$ T_0 = \frac{\hat{\beta}_1 - \beta_{1,0}}{\sqrt{\hat{\sigma}^2/S_{xx}}} \sim t_{n-2}$$

\vspace{0.05 in}
On rejette si $H_0$ versus $H_1 : \beta_1 > \beta_{1,0}$ si $t_0 > t_{\alpha, n-2}$.



%\vspace{0.1cm}
\textbf{Intervalles de confiance:}
\vspace{0.1cm}

{\renewcommand{\arraystretch}{2.4}
\begin{tabular}{|c|c|c|}
\hline Quantité \`a estimer & Estimateur & Intervalle de
confiance de niveau $1-\alpha$\\\hline

\hline $\beta_1$ & $\displaystyle \hat{\beta}_1 =\frac{
S_{xy}}{S_{xx}}$& $\hat{\beta}_1 \pm t_{\alpha/2,n-2}
\sqrt{\displaystyle{\frac{\hat{\sigma}^2}{S_{xx}}}}$ \\[.20cm]

\hline
%$E(Y|x_0)$ & $\widehat{E(Y|x_0)}=\hat y_0=\hat{\beta}_0 +\hat{\beta}_1 x_0$
Valeur moyenne de $Y$ en $x_0$ & $\hat y_0=\hat{\beta}_0 +\hat{\beta}_1 x_0$
& $\hat y_0 \pm t_{\alpha/2,n-2}\sqrt{\hat{\sigma}^2\left(\displaystyle{\frac{1}{n} + \frac{(x_0 -\overline{x})^2}{S_{xx}}}\right)}$ \\[.20cm]

\hline Observation future $y_0$ & $\hat y_0=\hat{\beta}_0 +
\hat{\beta}_1 x_0$ & $\hat y_0 \pm t_{\alpha/2,n-2}
\sqrt{\hat{\sigma}^2\left(\displaystyle{1 + \frac{1}{n} + \frac{(x_0 -
\overline{x})^2}{S_{xx}}}\right)}$ \\[.20cm]

\hline
\end{tabular}
}
\begin{figure}[p]
\makebox[\linewidth]{\includegraphics[scale = 0.9]{Student.pdf}}
\end{figure}
\begin{figure}[p]
\makebox[\linewidth]{\includegraphics[scale = 0.9]{Khicarre.pdf}}
\end{figure}

\end{document} 